\documentclass[aspectratio=169, 10pt]{beamer}
\usetheme{metropolis}

% --- Edinburgh colours ----------------------------------------
\definecolor{UoERed}{HTML}{7A2318}
\definecolor{UoEGold}{HTML}{B8860B}
\definecolor{CodeBG}{HTML}{F7F3ED}
\definecolor{CodeFrame}{HTML}{D5C9B1}

\setbeamercolor{palette primary}{bg=UoERed, fg=white}
\setbeamercolor{frametitle}{bg=UoERed, fg=white}
\setbeamercolor{title separator}{fg=UoEGold}
\setbeamercolor{progress bar}{fg=UoEGold, bg=UoEGold!20}
\setbeamercolor{alerted text}{fg=UoERed}
\setbeamercolor{block title}{bg=UoERed!15, fg=UoERed}
\setbeamercolor{block body}{bg=UoERed!5}

% --- Packages -------------------------------------------------
\usepackage[T1]{fontenc}
\usepackage{lmodern}
\usepackage{amsmath, amssymb, mathtools}
\usepackage{listings}
\usepackage{graphicx, booktabs}
\usepackage{tikz}
\usetikzlibrary{arrows.meta, positioning, calc, decorations.pathreplacing}
\usepackage{hyperref}
\hypersetup{colorlinks=true, linkcolor=UoERed, urlcolor=UoERed!70!black}
\setbeamertemplate{navigation symbols}{}

\lstset{
  language=Python,
  basicstyle=\ttfamily\scriptsize,
  keywordstyle=\color{UoERed}\bfseries,
  commentstyle=\color{gray}\itshape,
  stringstyle=\color{UoEGold},
  backgroundcolor=\color{CodeBG},
  frame=single,
  rulecolor=\color{CodeFrame},
  numbers=none, breaklines=true, showstringspaces=false, tabsize=4,
  xleftmargin=4pt, xrightmargin=4pt, aboveskip=6pt, belowskip=4pt,
}

\AtBeginSection[]{
  \begin{frame}
    \vfill\centering
    \usebeamerfont{title}\insertsectionhead\par
    \vfill
  \end{frame}
}

\title{Week 7 --- Physics-Informed Neural Networks}
\subtitle{Continuous-Time Economics \& the HJB Equation}
\author{Dr Juan Zurita}
\institute{Deep Learning for Macroeconomics\\
The University of Edinburgh~$\cdot$~School of Economics}
\date{}

\begin{document}
\maketitle

\begin{frame}{From Discrete to Continuous Time}
Many economic models are naturally continuous-time:

$$\rho\,V(x) = \max_c\!\left\{u(c) + V'(x)\,\mu(x,c) + \tfrac{1}{2}\,V''(x)\,\sigma^2(x)\right\}$$

\bigskip
This is the \textbf{Hamilton--Jacobi--Bellman} (HJB) equation.

\bigskip
Traditional approach: finite-difference grids. PINN approach: train a neural network to satisfy the PDE.
\end{frame}

\begin{frame}{The PINN Principle}
\begin{block}{Raissi, Perdikaris \& Karniadakis (2019)}
Parameterise the solution as $V(x) = \mathcal{N}_{\boldsymbol{\theta}}(x)$ and minimise the PDE residual.
\end{block}

\bigskip
Loss = $\underbrace{\|PDE\text{ residual}\|^2}_{\text{interior}}$ + $\underbrace{\|BC\text{ residual}\|^2}_{\text{boundary}}$

\bigskip
Derivatives $V'$, $V''$ are computed via \textbf{autodiff} --- no finite differences needed.
\end{frame}

\begin{frame}{Cake-Eating as an HJB}
$\rho\,V(W) = \max_c\left\{u(c) + V'(W)\cdot(-c)\right\}$

\bigskip
FOC: $u'(c) = V'(W)$ $\Rightarrow$ $c^* = (V'(W))^{-1/\gamma}$

\bigskip
Substitute back: $\rho\,V = u(c^*) - c^*\,V'$

\bigskip
\begin{block}{Boundary condition}
$V(0) = 0$ (no cake left $\Rightarrow$ no value).
\end{block}

\bigskip
The PINN learns $V(W)$ that satisfies both the HJB and the boundary.
\end{frame}

\begin{frame}{Soft vs Hard Boundary Conditions}
\textbf{Soft BC}: include boundary violations in the loss (weighted penalty).

\textbf{Hard BC}: embed the boundary in the architecture: $V(W) = W \cdot \mathcal{N}_{\boldsymbol{\theta}}(W)$.

\bigskip
Hard BCs guarantee satisfaction by construction. Soft BCs are more flexible but may violate boundaries during training.

\bigskip
For economics: hard BCs are preferred when boundary conditions have clear economic meaning.
\end{frame}

\begin{frame}{Summary}
\begin{itemize}
\item PINNs solve PDEs by training neural networks on residuals
\item Autodiff provides exact spatial derivatives --- no grids needed
\item The HJB equation is the natural target for continuous-time models
\item Hard boundary conditions embed economic constraints directly
\end{itemize}

\bigskip
\textbf{Next week:} heterogeneous agents --- where deep learning truly shines.
\end{frame}

\end{document}
